%% main.tex -- skeleton for utad-tfg.cls
%%
%% This mirrors "Template - inso.docx" page for page, so you can open the
%% two side by side and check. Replace the guidance text as you write; the
%% structure is the part to keep.
%%
%% Build:  latexmk -pdf main.tex
%%         (pdflatex, biber, pdflatex, pdflatex -- latexmk handles it)
%%
%% [spanish] switches Index / Image index / Glossary / Appendix / Keywords
%% to their Spanish names. [tnr] uses real Times New Roman (XeLaTeX only).

\documentclass[english]{utad-tfg}

\usepackage[style=apa,backend=biber]{biblatex}
\addbibresource{refs.bib}
% Let biblatex break a URL between letters, not only at punctuation; a
% long path segment otherwise runs into the margin.
\setcounter{biburlnumpenalty}{9000}
\setcounter{biburlucpenalty}{9000}
\setcounter{biburllcpenalty}{9000}

% ---- Cover ---------------------------------------------------------
\tfgtitle{Title}
\tfgmodality{conventional or linked to projects}
\tfgdegree{DEGREE AND SPECIALIZATION}
\tfgcall{DATE OF SUBMISSION (ORDINARY OR EXTRAORDINARY AND ACADEMIC YEAR)}
\tfgstudent{NAME AND SURNAME}
\tfgtutor{NAME AND SURNAME}
\tfgcotutor{NAME AND SURNAME}   % leave empty and the line disappears

\begin{document}

\tfgcover

% ---- p2: Acknowledgements ------------------------------------------
\begin{tfgacknowledgements}
(optional)
\end{tfgacknowledgements}

% ---- p3: style note, then the Index ---------------------------------
\textit{Each kind of heading in this template has a style attached to it
(title 1 to 5, plus one for long quotations). In LaTeX you get that for
free: use \texttt{\textbackslash chapter}, \texttt{\textbackslash section},
\texttt{\textbackslash subsection} and so on, and the index below rebuilds
itself on every compile -- there is nothing to refresh by hand.}

% The template follows its note with a second paragraph holding a support
% URL; this one keeps a second paragraph so the index still runs onto the
% next page at the same entry.
\url{https://www.overleaf.com/learn/latex/Table_of_contents}

\tfgindex

% ---- p5: Abstract and Resumen ---------------------------------------
% \tfgabstract is the summary in the document's language and comes first;
% \tfgaltabstract is the other one. With [spanish] the headings swap to
% Resumen / Abstract, so swap the two texts as well.
\begin{tfgabstractpage}
\tfgabstract{%
  Brief summary of the BFP in English. It is recommended to describe in a
  few words the topic, your research and the conclusions you have reached.%
}{four or five keywords}

\tfgaltabstract{%
  Breve resumen del TFG en espa\~nol. Se recomienda describir en pocas
  palabras la tem\'atica, el trabajo desarrollado y las conclusiones.%
}{cuatro o cinco palabras clave}
\end{tfgabstractpage}

% =====================================================================
\chapter{Introduction}

\section{Justification and context}
This section must explain the sense of this BFP within a specific field of
study (animation, game studies, computer graphics, renewable energy\ldots)
and what it aims to contribute. The discourse should make it clear why this
topic is relevant.

\section{Motivation}
This is the only place of the BFP where you can write in the first person.
Here you explain the personal reasons that led the author to write this BFP.

\section{Problem statement}
Here you must explain which is the problem or problems that your BFP is
going to address. You can use as many subtitles as necessary.

\section{Objectives}
They are usually divided into general objectives (the main focus of your
paper) and other specific objectives, complementary to the general one.

% =====================================================================
\chapter{State of the art}

In a BFP the author joins a conversation which has already started, and it
only makes sense to join a conversation if we keep up to date with
everything that has been already discussed or made. This section provides
that update -- what others have achieved before the author arrived and
which are the essential concepts for understanding the field. It is crucial
and is based on careful research.

\section{First chapter of your state of the art [Ex.: Psychology of terror,
Electromagnetic waves, etc.]}
The state of the art must be structured from the most general information
to the most specific.

You can devote as many titles and subtitles as you need. Remember to use the heading commands provided by this class.

% =====================================================================
\chapter{Methodology}

\begin{figure}[H]
  \centering
  % The template's own figures are 14.99 x 12.70 cm -- full text width.
  % Matching that is what makes the two figures fall on separate pages
  % here exactly as they do in the .docx. The frame's padding and rule are
  % taken off both dimensions so the box is the picture's exact size.
  \fbox{\parbox[c][\dimexpr12.7cm-2\fboxsep-2\fboxrule][c]{\dimexpr\textwidth-2\fboxsep-2\fboxrule}{\centering\itshape
    figure placeholder -- replace with \texttt{\textbackslash includegraphics}}}
  \caption{}
  \label{fig:one}
\end{figure}

This chapter explains how the knowledge presented in the state of the art
will be reflected in the development. If the development involves a
product, then you should explain why some particular features were chosen
instead of others, which were the criteria that determined that choice.

\begin{figure}[H]
  \centering
  \fbox{\parbox[c][\dimexpr12.7cm-2\fboxsep-2\fboxrule][c]{\dimexpr\textwidth-2\fboxsep-2\fboxrule}{\centering\itshape
    figure placeholder -- replace with \texttt{\textbackslash includegraphics}}}
  \caption{A caption for the second figure. \source{Author (2026)}}
  \label{fig:two}
\end{figure}

You can add as many titles as you need.

% The template skips 3.1 and labels this section 3.2 -- an off-by-one in
% the original that its own index repeats. Reproduced here so the two
% documents line up; delete this line to number it 3.1.
\setcounter{section}{1}
\section{Employed technologies}

% =====================================================================
\chapter{Development}

This chapter is intended for students to conduct their own exploration,
relating it to the state of the art described above.

Here you can describe the different phases of the work, relating them to
the topics covered in the state of the art (``this has been worked on in
this way, differently or similarly to what the theory says / other
practical works analysed in the theoretical framework have done'').

These are only examples. You can find many different types of BFP
development.

You can add as many titles as you need.

\section{}

% =====================================================================
\chapter{Conclusions}

This section can be read independently of the rest of the document. It
should summarise the main topics of your research. The order in which the
ideas are presented may differ slightly from that of the document, and may
be organised according to different criteria.

Please note that new information should never appear in the conclusions
(this is the purpose of the development section), and that the tone of the
writing must be academic -- it should not be included whether the work was
important or not personally for the student. As in the rest of the BFP
(excepting the motivation), the first-person singular must not be used.

\section{Discussion}
This section mentions what could not be covered during the development of
the BFP. This may include any problems that arose during the realization of
the BFP: someone did not respond to an interview request, a technical
problem occurred that prevented the completion of a certain section of the
BFP, etc.

\section{Further research}
This section presents further research lines of this BFP in the future.

It should include whether the present work can be replicated for other
examples, or to what directions it can be extended in the future.

% =====================================================================
\chapter{References}

\section{Bibliography}
It is presented in alphabetical order in APA format. Here it is generated
from \texttt{refs.bib} by \texttt{biblatex}, so the ordering and the
formatting are not yours to maintain \parencite{hevner2004,peffers2007,%
wieringa2014,kimball2013,loshin2008}.

\printbibliography[heading=none]

\section{Image index}
\tfgimageindex

% =====================================================================
\begin{tfgglossary}
The BFP is a document intended for someone who is not an expert in the
subject. Therefore, here can be listed all the specific vocabulary that may
be unfamiliar to someone outside the field of expertise (rig, game loop,
neural network) but that would hinder the discourse of the state of the art.

It should be listed in alphabetical order and you do not need to cite
references.

There may be BFP where this section is not necessary.
\end{tfgglossary}

\begin{tfgappendix}
Include here all cited material that takes up too much space to be included
in the body of the text: any graph, sketch, statistic, line of code,
reference text, survey, interview, or technical note that is not directly
related to the development and conclusions of the BFP but that you wish to
add as reference material.
\end{tfgappendix}

\end{document}
